\documentclass[../../../include/open-logic-section]{subfiles}

\begin{document}

\olfileid[id]{sth}{replacement}{extrinsic}
\olsection[Pertimbangan Ekstrinsik]{Pertimbangan Ekstrinsik mengenai Penggantian}

Kita mulai dengan mempertimbangkan upaya \emph{ekstrinsik} untuk membenarkan
Penggantian. Boolos mengusulkan upaya berikut.
\begin{quote}
  [\ldots] alasan untuk mengadopsi aksioma-aksioma penggantian cukup
  sederhana: aksioma-aksioma itu mempunyai banyak konsekuensi yang
  dikehendaki dan (tampaknya) tidak mempunyai konsekuensi yang tidak
  dikehendaki. Selain teorema-teorema mengenai konsepsi iteratif,
  konsekuensinya mencakup teori bilangan tak hingga yang memuaskan meskipun
  tidak ideal, dan suatu hasil yang sangat dikehendaki yang membenarkan
  definisi induktif pada relasi yang berfondasi baik.
  \citep[229]{Boolos1971}
\end{quote}
Inti gagasan Boolos ialah bahwa kita seharusnya membenarkan Penggantian
berdasarkan buahnya. Buah tertentu yang ia sebutkan merupakan hal-hal yang
telah kita bahas dalam beberapa bab terakhir. Penggantian memungkinkan kita
membuktikan bahwa ordinal von Neumann merupakan pengganti yang sangat baik
bagi gagasan tipe pengurutan baik (inilah ``teori bilangan tak hingga yang
memuaskan meskipun tidak ideal'' kita). Penggantian juga memungkinkan kita
mendefinisikan $V_\alpha$, menetapkan gagasan peringkat, dan membuktikan
Induksi-$\in$ (semua ini merupakan ``teorema-teorema mengenai konsepsi
iteratif'' kita). Terakhir, Penggantian memungkinkan kita membuktikan Teorema
Rekursi Transfinit (inilah ``definisi induktif pada relasi yang berfondasi
baik'').

Semua ini memang konsekuensi yang dikehendaki. Namun, apakah konsekuensi yang
dikehendaki itu cukup untuk \emph{membenarkan} Penggantian? \emph{Tidak}.
Atau setidaknya, tidak secara langsung.

Berikut suatu masalah sederhana. Walaupun kita telah menyatakan beberapa
konsekuensi Penggantian yang dikehendaki, banyak di antaranya dapat kita
peroleh melalui cara lain. Hal ini tidak dikenal seluas yang semestinya,
sehingga kita perlu berhenti sejenak untuk menjelaskan keadaannya.

Terdapat teori himpunan sederhana yang disebut Teori Tingkat, atau disingkat
$\LT$.\footnote{Versi-versi pertama $\LT$ diajukan oleh
\citet{Montague1965} dan \citet{Scott1974}; teori ini disederhanakan dan
diberi pembahasan sepanjang buku oleh \citet{Potter2004}; dan
\citet{ButtonLT1} baru-baru ini menyederhanakan $\LT$ lebih lanjut.}
Aksioma-aksioma $\LT$ hanyalah Ekstensionalitas, Pemisahan, dan klaim bahwa
setiap himpunan merupakan himpunan bagian dari suatu \emph{tingkat}; di sini
``tingkat'' didefinisikan secara cerdik sehingga tingkat-tingkat berperilaku
seperti sahabat kita, $V_\alpha$. Jadi, $\ZF$ membuktikan $\LT$; tetapi
$\LT$ \emph{jauh} lebih lemah daripada $\ZF$. Bahkan, $\LT$ tidak memberi
Anda Pasangan, Himpunan Kuasa, Ketakhinggaan, ataupun Penggantian. Misalkan
$\Zr$ merupakan hasil penambahan Ketakhinggaan dan Himpunan Kuasa kepada
$\LT$; ini juga memberikan Pasangan, sehingga $\Zr$ setidaknya sekuat $\Z$.
Namun, sebenarnya $\Zr$ benar-benar lebih kuat daripada $\Z$, karena teori
itu menambahkan klaim bahwa setiap himpunan mempunyai peringkat (karena itu
saya menyarankan agar teori tersebut disebut $\Zr$). Sesungguhnya, $\Zr$
memberikan: teori ordinal yang sepenuhnya memuaskan; hasil-hasil yang
menstratifikasi hierarki menjadi tahap-tahap yang terurut baik; bukti
Induksi-$\in$; dan suatu \emph{versi} Rekursi Transfinit.

Singkatnya: walaupun Boolos tidak mengetahui hal ini, semua konsekuensi yang
dikehendaki yang ia sebutkan dapat diperoleh \emph{tanpa} Penggantian; ia
hanya perlu menggunakan $\Zr$, bukan $\Z$.

(Dengan mempertimbangkan semua ini, mengapa kita mengikuti jalur
konvensional dengan mengajarkan $\ZF$ kepada Anda, alih-alih $\LT$ dan
$\Zr$? Ada dua alasan. Pertama: semata-mata karena alasan historis, memulai
dengan $\LT$ cukup tidak baku; kita ingin membekali Anda agar dapat membaca
pembahasan teori himpunan yang lebih baku. Kedua: ketika Anda siap memahami
$\LT$ dan $\Zr$, Anda dapat langsung membaca \citealt{Potter2004} dan
\citealt{ButtonLT1}.)

Tentu saja, karena $\Zr$ benar-benar lebih lemah daripada $\ZF$, terdapat
hasil-hasil yang dibuktikan $\ZF$ tetapi dibiarkan terbuka oleh $\Zr$. Jadi,
seseorang dapat mencoba membenarkan Penggantian atas dasar ekstrinsik dengan
menunjuk salah satu hasil tersebut. Namun, setelah Anda mengetahui cara
menggunakan $\Zr$, cukup sulit menemukan banyak contoh hal yang (a)
ditetapkan oleh Penggantian dan bukan oleh cara lain, serta (b) secara
intuitif benar. (Untuk pembahasan lebih lanjut, lihat
\citealt[\S13.2]{Potter2004}.)

Kesimpulannya sebagai berikut. Untuk memberikan pembenaran ekstrinsik yang
meyakinkan bagi Penggantian, kita perlu menemukan suatu hasil yang
\emph{tidak dapat} dicapai tanpa Penggantian. Dan itu bukan usaha yang mudah.

Mari kita pertimbangkan satu masalah lagi yang timbul bagi setiap upaya untuk
memberikan pembenaran yang semata-mata ekstrinsik bagi Penggantian. (Masalah
ini mungkin lebih mendasar daripada yang pertama.) Boolos bukan sekadar
menunjukkan bahwa Penggantian mempunyai banyak konsekuensi yang dikehendaki.
Ia juga menyatakan bahwa Penggantian ``(tampaknya) tidak mempunyai''
konsekuensi yang tidak dikehendaki. Namun, kualifikasi dalam tanda kurung,
``tampaknya,'' tentu sungguh menentukan.

Ingat bagaimana kita sampai di sini: Komprehensi Naif mengalami
inkonsistensi, dan kita menanggapi inkonsistensi ini dengan menerima
konsepsi kumulatif-iteratif tentang himpunan. Konsepsi ini disertai sebuah
kisah yang, kita harapkan, meyakinkan kita akan konsistensinya. Namun, jika
kita tidak dapat membenarkan Penggantian dari dalam kisah itu, maka kita
(sejauh ini) tidak mempunyai alasan untuk percaya bahwa $\ZF$ konsisten.
Atau lebih tepatnya: kita tidak mempunyai alasan untuk percaya bahwa $\ZF$
konsisten, selain fakta (yang mungkin sekadar kebetulan) bahwa belum ada
seorang pun yang menemukan kontradiksi \emph{sampai sekarang}. Tepat dalam
pengertian itu, komentar Boolos tampaknya dapat diringkas menjadi:
``(tampaknya) $\ZF$ konsisten''. Kita harus menuntut jaminan konsistensi yang
lebih kuat daripada ini.

Masalah ini akan memengaruhi setiap upaya yang \emph{murni} ekstrinsik untuk
membenarkan Penggantian, yakni setiap pembenaran yang dirumuskan semata-mata
dalam konsekuensi (yang diketahui) dari $\ZF$. Karena itu, kita akan mencari
pembenaran \emph{intrinsik} bagi Penggantian, yakni pembenaran yang
menyarankan bahwa kisah yang kita ceritakan mengenai himpunan entah bagaimana
``sudah'' mengikat kita pada Penggantian.

\end{document}
